% Page 026 — page imprimée 22
% Contenu seul : le préambule et la pagination sont gérés par meyrueis.tex

\sectiondate{Histoire de Samuel François}{1789-1795}

\begin{center}

\vspace{0.45em}

{\large\bfseries Pasteur de Meyrueis, nommé Administrateur\\
du Département de la Lozère}

\vspace{0.45em}

{\small\bfseries\itshape
D’après les comptes-rendus de la Société Populaire de Meyrueis\\
Archives Gévaudanaises : 1913. JB Belon P. 293 à 316}
\end{center}

\vspace{0.7em}

A cette époque Meyrueis avait bénéficié du ministère pastoral de Samuel François.
Écœuré sans doute, par les excès de la Révolution il avait été nommé chef du parti Girondin de
Lozère. Il soutenait un certain fédéralisme au point qu’il fut jeté en prison à Mende. Il réussit
à s’évader et les Montagnards se jetèrent à sa poursuite.

C’est ainsi qu’arriva à Meyrueis, le 6 germinal (27 mars 1794), le Montagnard Chateauneuf-
Randon qui, prenant la parole « avance, dit-il, avec regret, qu’il est convaincu qu’il existe
dans la commune des citoyens instruits du lieu de la retraite de Samuel François ; il démontre
avec fermeté le crime de la récelation (sic) et propose de faire une adresse au peuple pour
l’éclairer de son devoir. »

A ces mots la Société par un mouvement spontané se lève et jure de révéler tout ce qu’elle
pourra découvrir de relatif à la retraite de « cet homme pervers qui a tenté, mais en vain,
d’avilir la représentation nationale et d’égarer le peuple sur ses véritables intérêts ».

Samuel François reste introuvable ; Chateauneuf-Randon décide alors de faire incarcérer sa
femme, mère de plusieurs enfants. Son sort émeut quelques anciens paroissiens de son mari à
Meyrueis, et lors d’une séance du Comité, le 24 vendémiaire an III (16 Octobre 1794) l’un
d’entre eux prend la parole et s’exprime en ces termes :

\begin{quote}
\itshape
« Citoyens, je viens intéresser votre justice, votre humanité en faveur de la citoyenne
François, dont le civisme, depuis 1789, ne s’est jamais démenti. Vous avez été les témoins de
sa conduite pendant la Révolution. Vous connaissez ses malheurs et la pureté de ses
intentions ; vous êtes instruits des motifs qui l’ont privée de sa liberté.

Vous savez qu’elle est en arrestation par ordre du représentant du peuple Châteauneuf-
Randon, depuis plus de six mois. Vous n’ignorez pas que l’évasion de son mari fut la seule
cause de cette mesure de sûreté générale. Vous savez encore que cette citoyenne a une
famille qui ne vit que de son travail et que son absence a réduite à la dernière misère. Ces
motifs de justice et d’humanité basent une réclamation qui tend à ce que la Société se
prononce en faveur d’une citoyenne intéressante par ses malheurs et son civisme, et qu’elle
sollicite sa mise en liberté auprès du comité de Sûreté Générale de la Convention. Plusieurs
membres parlent de la citoyenne François et tous manifestent le même sentiment à son égard.

L’assemblée arrête à l’unanimité que le comité de Sûreté Générale de la Convention sera
invité d’ordonner la mise en liberté de la citoyenne Samuel François de Meyrueis, détenue
dans la maison de réclusion de Mende par ordre du représentant du peuple Châteauneuf-
Randon. »
\end{quote}

% Page 027 — page imprimée 23
% Contenu seul : le préambule et la pagination sont gérés par meyrueis.tex

\textbf{Il avait donc fallu 6 mois pour qu’un membre du Comité Populaire se soucie du sort de
la femme du pasteur Samuel François, et ose, suivi par tous les membres du Comité
Populaire de Meyrueis, se dresser contre la mesure unique prise par Châteauneuf-
Randon. !!!}

Quelques mois plus tard, les poursuites contre Samuel François sont levées. Il peut reparaître
en plein jour, et regagne Meyrueis, sa famille et les séances de la Société Populaire. Voilà
comment est relaté son retour (séance du 21 ventôse An III (11 mars 1795).

\begin{quote}
\itshape
«\ldots\ Un sujet intéressant devait fixer dans cette séance l’attention et l’attendrissement de
chaque citoyen : c’était le récit de la persécution affreuse à laquelle a été en butte toute une
année, l’honnête citoyen dont tout le crime était d’avoir osé défendre les droits du peuple. Il a
demandé d’être réintégré dans la Société ; il a même produit des preuves de son civisme
reconnu pendant sa fuite ; mais l’assemblée assez pénétrée depuis longtemps des vertus
énergiques du citoyen Samuel François, l’a admis comme par l’acclamation et l’accolade
fraternelle qu’il a reçue des membres du bureau, lui est un sûr garant de la joie que ressentait
la Société du recouvrement d’un membre, dont l’absence n’avait pu, que se faire sentir
vivement dans son sein et dans l’administration du Département, dont il était membre. »
\end{quote}

\textbf{En mars 1794 personne n’avait osé, au moins en public, prendre la défense de Samuel
François alors qu’il venait de s’évader : ce n’était qu’un scélérat, mais un an plus tard,
en Mars 1795, les membres de la Société se jettent à son cou comme si de rien n’était.}

Il est vrai que \textbf{Châteauneuf-Randon} (\textit{de son vrai nom Alexandre Paul Guérin du
Tournel, marquis de Joyeuse, Comte de Châteauneuf-Randon}) et ses méthodes violentes
ne sévissait plus en Lozère. (Il lui fallait bien faire la preuve, lui, l’aristocrate, qu’il était un
vrai montagnard -- révolutionnaire). Il avait été \textbf{nommé auprès de l’armée des Pyrénées}.

\vfill

\begin{center}
\small
Extraits des comptes-rendus de la Société Populaire de Meyrueis,\\
édités par la Société d’Agriculture, Industrie, Sciences et\\
Arts de la Lozère, 1913. Archives Gévaudanaises\\
Mars 2006
\end{center}
